\usepackage{xeCJK}
\usepackage{geometry}
\geometry{left=2cm,right=2cm,top=2.5cm,bottom=2.5cm}
\setlength{\headheight}{12.7pt}

\usepackage{titlesec}
\titleformat{\section}
  {\normalfont\large\bfseries}
  {\thesection}
  {1em}
  {}
\titlespacing*{\section}
  {0pt}
  {3.5ex plus 1ex minus .2ex}
  {2.3ex plus .2ex}

\usepackage{amsmath}
\usepackage{amssymb}
\usepackage{amsthm}
\usepackage{enumitem}
\setlist[enumerate]{itemsep=0pt,topsep=0pt,parsep=0pt,leftmargin=0pt,itemindent=2.5em}
\setlist[itemize]{itemsep=0pt,topsep=0pt,parsep=0pt,leftmargin=0pt,itemindent=2.5em}
\usepackage{fancyhdr}
\pagestyle{fancy}
\fancyhf{}
\usepackage{graphicx}
\usepackage{hyperref}
\usepackage{indentfirst}
\usepackage{lmodern}


\begin{document}
\setlength{\abovedisplayskip}{1pt}
\setlength{\belowdisplayskip}{1pt}

\section{传递集}

\hypertarget{sec:定义1.1}{}
\noindent\textbf{定义1.1 (传递集)}

对任意的集合$T$, 如果
\[
    \forall x\,\forall t:\quad x\in t\land t\in T\to x\in T,
\]
就称$T$是\textbf{传递集}.

\vspace{10pt}

传递集有如下的等价定义.

\vspace{10pt}

\hypertarget{sec:定理1.1}{}
\noindent\textbf{定理1.1}

任给集合$T$, 以下等价:
\begin{enumerate}
    \item $T$是传递集,
    \item $\forall t\in T:t\subset T$,
    \item $\forall t\in T:t\in\mathcal{P}(T)$,
    \item $T\subset\mathcal{P}(T)$,
    \item $\displaystyle\bigcup T\subset T$.
\end{enumerate}

\noindent{\kaishu 证明.}

显然前4条互相等价, 另外4.与5.也等价, 因为
\vspace{3pt}
\begin{enumerate}
    \item 如果$T\subset\mathcal{P}(T)$, 那么
    $\displaystyle\bigcup T\subset\bigcup\mathcal{P}(T)=T$.
    \vspace{3pt}
    \item 如果$\displaystyle\bigcup T\subset T$, 那么
    $\displaystyle T\subset\mathcal{P}\left(\bigcup T\right)
    \subset\mathcal{P}(T)$. \hfill $\qedsymbol$
\end{enumerate}

\vspace{10pt}

\hypertarget{sec:定理1.2}{}
\noindent\textbf{定理1.2}

\begin{enumerate}
    \item 任给传递集$T$, $T^+$也是传递集,
    \item 任给集合$A$, 如果$A$的元素都是传递集, 那么$\displaystyle\bigcup A$是传递集,
    \item 任给集合$A$, 如果$A$的元素都是传递集, 那么$\displaystyle\bigcap A$是传递集.
\end{enumerate}

\noindent{\kaishu 证明.}

\begin{enumerate}
    \item 对任意的$t\in T^+$, 有$t\in T$或$t=T$,
    所以$t\subset T$, 从而$t\subset T^+$.
    \vspace{3pt}
    \item 对任意的$x\in t\in\displaystyle\bigcup A$, 存在$T\in A$使得$t\in T$,
    于是$x\in T$, 从而$\displaystyle x\in\bigcup A$.
    \vspace{3pt}
    \item 对任意的$x\in t\in\displaystyle\bigcap A$, 对任意的$T\in A$都有$t\in T$,
    于是$x\in T$, 从而$\displaystyle x\in\bigcap A$. \hfill $\qedsymbol$
\end{enumerate}

\end{document}